\section{Discussion and Future Work}\label{sec:discussion}
One of the primary challenges we encountered was the application's performance. While effective in generating adequate quality summaries, the BART model introduced latency issues that impacted the user experience. This raises a problem with the app's responsiveness. Exploring more efficient models or optimizing the existing implementation might help strike a better balance.

One promising direction is the exploration of the Flan-T5 model for topic modeling. Flan-T5, known for its strong performance in various natural language processing tasks, could potentially offer more accurate and contextually relevant topic extraction, even for shorter texts. Fine-tuning Flan-T5 specifically for the task of topic modeling might address some of the limitations we encountered with BERTopic, particularly in generating meaningful topics from minimal input.

For the summarization task, several models could be considered to improve both the quality and speed of the generated summaries. The Flan-T5 model, when fine-tuned for summarization, could offer a powerful alternative to the BART model, potentially delivering faster processing times without compromising the quality of the summaries. Additionally, Pegasus, a model pre-trained specifically for abstractive summarization, could be another viable option. Pegasus is designed to handle summarization tasks effectively and might yield more concise and coherent summaries.

Moreover, the use of large language models (LLMs), such as GPT-4 or similar advanced models, could further enhance the summarization capabilities. These models have demonstrated exceptional performance in generating human-like text and could provide summaries that are not only accurate but also nuanced and contextually rich. However, integrating LLMs would require careful consideration of computational resources and cost, as well as strategies to mitigate potential biases in the generated content.\citep{shakil2024evaluating}

The limitations of BERTopic in processing shorter texts highlighted a significant challenge in natural language processing: extracting meaningful information from minimal input. Our fallback solution of using predefined topics was a practical response, but it also introduced a degree of inflexibility. In future work, incorporating techniques that can adaptively enhance the context or dynamically adjust the topic modeling process for shorter texts could improve the robustness of the application.

For the user interface, it can be more intuitive and responsive. Future work could explore more advanced UX/UI features, such as more sophisticated visualization tools, or even personalized summarization based on user preferences.

As the app evolves, scalability will become an increasingly important factor. Expanding the app to handle larger volumes of data or integrating it with more platforms will require careful infrastructure and software architecture planning. Additionally, the potential applications of the app extend beyond journalism; it could be adapted for use in various industries, such as market research, public relations, or any field that requires real-time analysis of social media content.
